\documentclass[conference]{IEEEtran}
\usepackage[T1]{fontenc}
\usepackage[utf8]{inputenc}
\usepackage[spanish,es-tabla]{babel}
\usepackage{amsmath}
\usepackage{amssymb}
\usepackage{booktabs}
\usepackage{array}
\usepackage{microtype}
\usepackage{tikz}
\usepackage{hyperref}
\usepackage{enumitem}
\usetikzlibrary{arrows.meta,positioning,shapes.geometric,calc}
\hypersetup{hidelinks}

\title{Diseño e Implementación de Nodos Worker Distribuidos para Procesamiento Paralelo de Imágenes con gRPC y \texttt{asyncio}}

\author{
\IEEEauthorblockN{Santos Enmanuel Manosalva Aceros}
\IEEEauthorblockA{
\textit{Universidad Pontificia Bolivariana}\\
Bucaramanga, Colombia
}
\and
\IEEEauthorblockN{Christian Mauricio Rodriguez Curubo}
\IEEEauthorblockA{
\textit{Universidad Pontificia Bolivariana}\\
Bucaramanga, Colombia
}
}

\begin{document}
\maketitle

\begin{abstract}
Este documento describe el diseño y la implementación de un subsistema de nodos \emph{worker} para procesamiento distribuido de imágenes en paralelo. La solución fue desarrollada en Python con \texttt{asyncio} y \texttt{grpc.aio}, y se integra con un servidor principal externo escrito en Java que concentra autenticación, base de datos y reglas de negocio. El prototipo final incluye tres workers homogéneos, contratos gRPC de negocio y control, cola de prioridad local, scheduling interno, control de admisión, persistencia durable, almacenamiento compartido, OCR e inferencia desacoplados, métricas, healthchecks, mTLS y despliegue contenedorizado. El resultado es un subsistema ejecutable, validado por pruebas automatizadas y listo para ser consumido por el servidor principal de la aplicación.
\end{abstract}

\begin{IEEEkeywords}
sistemas distribuidos, gRPC, procesamiento de imágenes, Python, asyncio, workers, scheduling
\end{IEEEkeywords}

\section{Introducción}
En sistemas con procesamiento intensivo de imágenes resulta conveniente separar la lógica de negocio del plano de ejecución. En este proyecto, el servidor principal de la aplicación permanece fuera del alcance del repositorio y se asume implementado en Java. Dicho servidor administra autenticación, persistencia relacional y coordinación global. El objetivo del trabajo fue construir únicamente el subsistema de ejecución: un clúster de workers capaz de recibir solicitudes por gRPC, priorizarlas, admitirlas según recursos y ejecutar transformaciones de imagen en paralelo.

La solución final no reemplaza al backend principal. Por el contrario, expone interfaces gRPC para que dicho backend pueda conectarse a los workers, enviar tareas, consultar estado y recibir callbacks opcionales de progreso y resultado.

\section{Alcance}
El alcance implementado cubre:
\begin{itemize}[leftmargin=*]
\item tres nodos worker con el mismo runtime;
\item contrato gRPC de negocio y contrato gRPC de control;
\item cola de prioridad local por nodo;
\item scheduler híbrido con prioridad, aging, deadline y costo estimado;
\item control de admisión, backpressure y retries;
\item persistencia durable de pendientes, resultados y eventos;
\item OCR e inferencia mediante adaptadores configurables;
\item despliegue con Docker, MinIO, Prometheus y Grafana.
\end{itemize}

Quedaron fuera de alcance la autenticación de usuario final, la base de datos corporativa, la coordinación global del clúster y la lógica de negocio del sistema, ya que dichas responsabilidades pertenecen al servidor principal externo.

\section{Arquitectura general}
La Fig.~\ref{fig:architecture} resume la arquitectura final del subsistema. El servidor principal Java se conecta a los workers por gRPC. Los workers son los únicos componentes ejecutables de este repositorio. MinIO resuelve almacenamiento compartido y Prometheus/Grafana proveen observabilidad.

\begin{figure*}[t]
\centering
\begin{tikzpicture}[
  font=\small,
  app/.style={draw, rounded corners=4pt, thick, minimum width=4.0cm, minimum height=1.1cm, align=center, fill=blue!10},
  worker/.style={draw, rounded corners=4pt, thick, minimum width=3.1cm, minimum height=1.15cm, align=center, fill=green!10},
  sub/.style={draw, rounded corners=3pt, thick, minimum width=3.1cm, minimum height=0.7cm, align=center, fill=white},
  store/.style={draw, rounded corners=4pt, thick, minimum width=4.2cm, minimum height=1cm, align=center, fill=orange!10},
  legend/.style={draw, rounded corners=3pt, thick, fill=white, align=left, inner sep=5pt},
  control/.style={-{Latex[length=2.5mm]}, ultra thick, blue!70!black},
  callback/.style={-{Latex[length=2.5mm]}, ultra thick, green!50!black},
  sharedarrow/.style={-{Latex[length=2.3mm]}, thick, gray!70}
]
\node[app] (app) at (0,4.6) {Servidor principal Java\\auth + BD + lógica de negocio};
\node[worker] (w1) at (-5.2,2.45) {Worker 1\\\texttt{127.0.0.1:50051}};
\node[worker] (w2) at (0,2.45) {Worker 2\\\texttt{127.0.0.1:50061}};
\node[worker] (w3) at (5.2,2.45) {Worker 3\\\texttt{127.0.0.1:50071}};

\node[sub] (w1q) at (-5.2,1.15) {Cola + scheduler + executor};
\node[sub] (w2q) at (0,1.15) {Cola + scheduler + executor};
\node[sub] (w3q) at (5.2,1.15) {Cola + scheduler + executor};

\node[store] (storage) at (-2.6,-0.85) {Storage compartido\\filesystem, \texttt{file://}, \texttt{s3://}, MinIO};
\node[store] (obs) at (3.7,-0.85) {Observabilidad\\Prometheus, Grafana, healthchecks};
\node[legend, anchor=west] (lg) at (-7.05,-2.15) {
\tikz{\draw[control] (0,0) -- (0.9,0);} RPCs desde el servidor principal: \texttt{ImageNodeService} y \texttt{WorkerControlService}\\
\tikz{\draw[callback] (0,0) -- (0.9,0);} Callbacks opcionales del worker: \texttt{ReportProgress}, \texttt{ReportResult}, \texttt{Heartbeat}\\
\tikz{\draw[sharedarrow] (0,0) -- (0.9,0);} Persistencia, resultados, métricas y healthchecks
};

\draw[control] ($(app.south)+(-0.85,0)$) to[out=-160,in=90] ($(w1.north)+(0.35,0)$);
\draw[control] ($(app.south)+(0.0,0)$) -- ($(w2.north)+(0.0,0)$);
\draw[control] ($(app.south)+(0.85,0)$) to[out=-20,in=90] ($(w3.north)+(-0.35,0)$);

\draw[callback] ($(w1.north)+(-0.35,0)$) to[out=90,in=-160] ($(app.south)+(-1.65,0)$);
\draw[callback] ($(w2.north)+(0.0,0)$) -- ($(app.south)+(0.0,0)$);
\draw[callback] ($(w3.north)+(0.35,0)$) to[out=90,in=-20] ($(app.south)+(1.65,0)$);

\draw[sharedarrow] (w1q.south) to[out=-80,in=145] (storage.north west);
\draw[sharedarrow] (w2q.south) -- ($(storage.north)+(0.8,0)$);
\draw[sharedarrow] (w3q.south) to[out=-100,in=35] (storage.north east);

\draw[sharedarrow] (w1.east) to[out=-18,in=160] (obs.west);
\draw[sharedarrow] (w2.east) -- (obs.west);
\draw[sharedarrow] (w3.south) -- ($(obs.north)+(1.05,0)$);
\end{tikzpicture}
\caption{Arquitectura final del subsistema. Azul: llamadas gRPC desde el servidor principal hacia los workers. Verde: callbacks opcionales del worker hacia el servidor principal. Gris: storage compartido, métricas y salud.}
\label{fig:architecture}
\end{figure*}

Cada worker contiene los siguientes subsistemas:
\begin{itemize}[leftmargin=*]
\item \textbf{servicio gRPC de negocio}: expone \texttt{ImageNodeService};
\item \textbf{servicio gRPC de control}: expone \texttt{WorkerControlService};
\item \textbf{priority queue}: implementada con \texttt{heapq};
\item \textbf{scheduler local}: combina prioridad, aging, deadline y costo estimado;
\item \textbf{resource manager}: evalúa CPU, memoria, I/O y concurrencia activa;
\item \textbf{execution manager}: despacha transformaciones sobre \texttt{ThreadPoolExecutor} y procesos dedicados cuando se requiere aislamiento;
\item \textbf{reporter opcional}: mantiene un canal gRPC hacia el servidor principal solo si se configura un endpoint de callbacks;
\item \textbf{state store}: preserva resultados terminales, pendientes y eventos pendientes de entrega.
\end{itemize}

\section{Contratos gRPC y comunicación}
Los contratos gRPC quedaron divididos en dos archivos.

\subsection{Servicio de negocio}
El archivo \texttt{proto/imagenode.proto} define \texttt{ImageNodeService}. Sus RPCs principales son:
\begin{itemize}[leftmargin=*]
\item \texttt{ProcessToPath};
\item \texttt{ProcessToData};
\item \texttt{UploadLargeImage};
\item \texttt{StreamBatchProcess};
\item \texttt{ProcessBatch};
\item \texttt{HealthCheck};
\item \texttt{GetMetrics}.
\end{itemize}

Este contrato es el más natural para el servidor principal Java, porque abstrae el procesamiento y devuelve bytes o rutas resultantes sin exponer la estructura interna de la cola.

\subsection{Servicio de control}
El archivo \texttt{proto/worker\_node.proto} define \texttt{WorkerControlService}. Sus RPCs principales son:
\begin{itemize}[leftmargin=*]
\item \texttt{SubmitTask};
\item \texttt{GetNodeStatus};
\item \texttt{CancelTask};
\item \texttt{DrainNode};
\item \texttt{ShutdownNode}.
\end{itemize}

Este contrato permite al servidor principal controlar la cola, consultar capacidad o drenar nodos antes de mantenimiento.

\subsection{Callbacks opcionales}
El mismo archivo \texttt{proto/worker\_node.proto} define \texttt{CoordinatorCallbackService}. Aunque el nombre histórico se conserva en el proto, en la arquitectura final representa el endpoint de callbacks del servidor principal. Sus métodos son:
\begin{itemize}[leftmargin=*]
\item \texttt{ReportProgress};
\item \texttt{ReportResult};
\item \texttt{Heartbeat}.
\end{itemize}

Si el worker no recibe un \texttt{WORKER\_COORDINATOR\_TARGET}, este canal se desactiva y el nodo sigue operando de forma standalone.

\section{Cola de prioridad y scheduling local}
Cada worker mantiene una cola de prioridad local sobre \texttt{heapq}. El score implementado combina prioridad estática, urgencia por deadline, envejecimiento, costo estimado, reintentos y tamaño de imagen:
\begin{equation}
\label{eq:score}
\text{score} = w_p P + w_d D + w_a A - w_c C - w_r R - w_s S
\end{equation}
donde $P$ es la prioridad normalizada, $D$ la urgencia relativa al deadline, $A$ el aging acumulado, $C$ el costo estimado, $R$ el número de intentos previos y $S$ el peso relativo de la imagen.

La política final es híbrida:
\begin{itemize}[leftmargin=*]
\item la prioridad acelera trabajos críticos;
\item el aging evita starvation;
\item el deadline incrementa el score cuando una tarea se acerca a su vencimiento;
\item el costo estimado favorece una variante de SJF sin sacrificar prioridad.
\end{itemize}

El estimador de costo se actualiza con promedio exponencial:
\begin{equation}
\tau_{n+1} = \alpha t_n + (1-\alpha)\tau_n
\end{equation}
donde $t_n$ es el tiempo real de servicio observado y $\tau_n$ la predicción anterior.

\section{Admisión, recursos y concurrencia}
Antes de ejecutar una tarea, el worker evalúa si existe capacidad disponible en CPU, memoria, I/O y GPU opcional. La capacidad efectiva se aproxima por:
\begin{equation}
\label{eq:capacity}
C_{\text{ef}} = \min(f_{\text{cpu}}, f_{\text{ram}}, f_{\text{gpu}}, f_{\text{io}})\cdot C_{\text{base}}
\end{equation}

Además, el worker aplica límites duros de entrada:
\begin{itemize}[leftmargin=*]
\item tamaño máximo por request;
\item tamaño máximo por lote;
\item número máximo de filtros;
\item dimensiones máximas y píxeles máximos por imagen.
\end{itemize}

En cuanto a concurrencia, el diseño se divide así:
\begin{itemize}[leftmargin=*]
\item \texttt{asyncio} para networking, control gRPC, scheduler, health y reporter;
\item \texttt{heapq} para la cola priorizada;
\item \texttt{ThreadPoolExecutor} para las transformaciones integradas;
\item procesos dedicados solo para aislamiento explícito o extensiones no cooperativas.
\end{itemize}

\section{Persistencia, errores y recuperación}
La tolerancia a fallos se resolvió en tres niveles:
\begin{enumerate}[leftmargin=*]
\item \textbf{deduplicación}: resultados terminales persistidos por \texttt{idempotency\_key};
\item \textbf{restauración de pendientes}: tareas en cola o programadas para retry se recargan al reinicio;
\item \textbf{spool durable de eventos}: progreso y resultados se almacenan hasta recibir \texttt{ack} del servidor principal.
\end{enumerate}

Si el endpoint de callbacks se desconecta, el worker no pierde información: guarda el evento y reintenta la entrega con backoff exponencial y \emph{replay} al restablecer el canal. Si no existe endpoint configurado, el nodo funciona sin reporter.

\section{Operaciones soportadas}
Las transformaciones integradas del worker son:
\begin{itemize}[leftmargin=*]
\item escala de grises;
\item \texttt{resize};
\item \texttt{crop};
\item \texttt{rotate};
\item \texttt{flip};
\item \texttt{blur};
\item \texttt{sharpen};
\item brillo/contraste;
\item watermark de texto;
\item conversión de formato JPG, PNG, TIF, WEBP, BMP, GIF estático e ICO.
\end{itemize}

Adicionalmente, el worker integra dos extensiones por adaptador:
\begin{itemize}[leftmargin=*]
\item \textbf{OCR}: ejecutable configurable por \texttt{WORKER\_OCR\_COMMAND};
\item \textbf{inferencia}: ejecutable configurable por \texttt{WORKER\_INFERENCE\_COMMAND}.
\end{itemize}

\section{Observabilidad y despliegue}
El subsistema expone:
\begin{itemize}[leftmargin=*]
\item métricas Prometheus por worker;
\item \texttt{/livez} y \texttt{/readyz};
\item logs estructurados;
\item tracing distribuido opcional con OpenTelemetry;
\item mTLS opcional en ejecución local y activado en el stack Docker;
\item despliegue local con \texttt{docker compose};
\item tres instancias worker, MinIO, Prometheus y Grafana.
\end{itemize}

El estado de \texttt{/readyz} considera cola, memoria disponible y conectividad con el endpoint de callbacks solo si este fue configurado. En modo standalone, el worker puede quedar listo sin servidor principal embebido.

\section{Validación}
El repositorio fue validado con pruebas unitarias e integración sobre:
\begin{itemize}[leftmargin=*]
\item procesamiento de imágenes;
\item persistencia de estado y spool;
\item recuperación de pendientes;
\item cancelación;
\item contratos gRPC de negocio y control;
\item operación del worker con y sin endpoint externo de callbacks.
\end{itemize}

La ejecución final de la suite arrojó veintidós pruebas aprobadas. También se validó una demo extremo a extremo directamente contra un worker y se verificó el despliegue contenedorizado con tres nodos.

\section{Conclusiones}
Se implementó exitosamente un subsistema de workers distribuidos para procesamiento de imágenes con Python, \texttt{asyncio} y gRPC. El resultado no es un backend completo sino un plano de ejecución desacoplado, con tres workers funcionales, contratos gRPC bien definidos, control de admisión, persistencia durable, cancelación, observabilidad, OCR concreto con Tesseract, inferencia concreta y despliegue seguro en contenedores. La arquitectura lograda separa con claridad las responsabilidades: el servidor principal Java coordina y decide; los workers Python priorizan, admiten, ejecutan, persisten y reportan.

\end{document}
